\chapter{Методы хранения}

Необходимым условием для работы с временными рядами является организация хранения отслеживаемых данных.

В силу особенностей временных рядов и их обработки обозначим ключевые требования к средствам их хранения:
\begin{itemize}
    \item приоритетность скорости записи данных перед скоростью их обработки, так как операции записи преобладают над операциями чтения~\cite{myref7};
    \item упорядоченность хранимых данных по времени, так как данные записываются, считываются и обновляются не случайным образом, а во временном порядке~\cite{myref7};
    \item расширяемость и вместительность хранилища данных;
    \item гарантия сохранности данных. 
\end{itemize}

В данной работе в качестве методов хранения выделим и рассмотрим следующие подходы:

\begin{itemize}
    \item (использование баз данных, специализированных для временных рядов -- OpenTSDB);
    \item использование нереляционных БД -- Cassandra;
    \item (использование файлов).
\end{itemize}

\section{Хранение временных рядов в Cassandra}

Apache Cassandra --- базы данных, которая является:

\begin{itemize}
    \item \textit{распределенной} --- обеспечивает работу на нескольких машинах, в дальнейшем называемых узлами, которые пользователем представляются как единое целое~\cite{myref3};
    \item \textit{децентрализованной} --- все узлы являются равноправными~\cite{myref3};
    \item \textit{эластично масштабируемой} --- обеспечивает возможность добавлять и убирать узлы, не прерывая работы системы.
Данное свойство является специальным для горизонтальной масштабируемости --- механического увеличения пропускной способности оборудования путум добавления памяти в уже имеющуюся машину~\cite{myref3};
    \item \textit{отказоустойчивой} --- предполагает замену повреждунных узлов без простоя, за счут их репликации на несколько узлов~\cite{myref3}
(репликация представляет собой набор технологий копирования и распространения данных и объектов между узлами, а также включает синхронизацию этих узлов для поддержания их согласованности
(согласованность означает, что операция чтения всегда возвращает последнее записанное значение)~\cite{myref3, myref4});
    \item \textit{допускающей настройку согласованности} --- существует возможность настроить количество узлов, которые должны подтвердить завершение операции обновления данных~\cite{myref3}.
\end{itemize}
Основными структурами данных в Cassandra являются:
\begin{itemize}
\item \textit{столбец} --- пара имя-значение;
\item \textit{строка} --- контейнер для упорядоченной коллекции столбцов на которые можно сослаться по первичному ключу;
\item \textit{таблица} --- контейнер для упорядоченной коллекции строк;
\item \textit{пространство ключей} --- контейнер таблиц;
\item \textit{кластер} --- контейнер пространств ключей, расположенных в одном или нескольких узлах~\cite{myref3}.
\end{itemize}

Таким образом, столбец --- самая элементарная структурная единица в модели данных Cassandra, и ему присущи следующие характеристики:
имя; значение; тип; временная метка, генерируемая при каждом обновлении столбца; время жизни, указывающее срок хранения значения столбца~\cite{myref3}.

Уточнив основные понятия базы данных, рассмотрим особенности Cassandra применительно к хранению временных рядов.

Модель данных организует хранение временного ряда, где
\textit{ключ для строки} --- имя для нашего измерения,
\textit{имя столбца} --- время измерения,
\textit{значение столбца} --- измеренное значение~\cite{myref5}.
Данная структура позволяет выбрать как индивидуальное значение, так и набор значений для заданного интервала какого-либо измерения~\cite{myref5}.

Кроме того, для каждого столбца хранится время жизни, позволяющее автоматизировать процесс удаления устаревших данных, а также временная метка, которая используется с целью:

\begin{itemize}
\item разрешения конфликтов при изменении одного и того же значения со стороны нескольких клиентов~\cite{myref3},
\item организации автоматической сортировки данных по времени~\cite{myref5}. 
\end{itemize}

В случае обработки временных рядов скорость поступления данных имеет большее значение, нежели скорость их изменения,
потому что информация однородна и фактически обрабатывается только небольшая часть из всего объума хранимой информации~\cite{myref5}.
Данные особенности позволяют выделить Cassandra, как нереляционную БД (скорость добавления данных в NoSQL системах выше, чем в SQL системах)~\cite{myref5}.
Cassandra является распределунной и децентрализованной, вследствие этого можно выделить следующие еу аспекты:
\begin{itemize}
\item отказоустойчивость --- гарантия сохранности информации за счут репликации~\cite{myref6};
\item масштабируемость --- увеличение объума хранимых данных возможно посредством добавления узла в кластер~\cite{myref6};
\item слабая согласованность данных --- отсутствие контроля актуальности читаемой информации~\cite{myref3}.
Данная особенность может быть скорректирована настройкой согласованности. 
\end{itemize}

\section{Выводы}
Сравнительный анализ методов хранения временных рядов будет
проведён во время групповой работы.

\begin{table}[h!]
    \caption{Анализ системы хранения временных рядов Cassandra}
    \label{tab_mem}
    \centering
    \begin{tabular}{||p{7cm}|p{7cm}||}
        \hline
        {параметр / система хранения} & Cassandra\\
        \hline \hline
        {модель данных} & строковая \\
        \hline
        {автоматическое удаление устаревших данных} & есть \\
        \hline
        {автоматическая сортировка данных по времени} & есть \\
        \hline
        {сравнение скорости поступления и обработки данных} &  {скорость поступления данных выше} \\
        \hline
        {отказоустойчивость} & есть \\
        \hline
        {масштабируемость} & {гибкая масштабируемость} \\
        \hline \hline
    \end{tabular}
\end{table}

Выполнив анализ базы данных Cassandra, представляется возможным
сделать вывод о том, что она удовлетворяет основным требованиям к
хранилищу временных рядов.